\subsection{From constraints to MIPS32: conformance and determinism}
\label{sec:verification}

The preceding subsections show, under the hypotheses of Theorem~\ref{thm:composition}, that an accepting proof implies a trace satisfying every constraint, and Theorem~\ref{thm:unique} makes that trace unique given chip determinism. This subsection argues that such a trace is the MIPS32 execution, in two parts that no soundness theorem can supply: the second is proved, the first evidenced. The constraints must compute the instruction set's function: a chip that computes the wrong one is sound for the wrong machine (\S\ref{sec:verification-conformance}). And they must admit only that function: a chip with two outputs for one input lets a cheating prover choose, which honest traces never expose (\S\ref{sec:verification-det}). Conformance suites carried through proving and an independent executable model provide the evidence; the proof is in \Lean{}, chip by chip.

\paragraph{Correctness requirements.} Soundness rests only on the verifier and what it reads: every chip's constraint description, the preprocessed tables and program commitment in the verifying key, the recursion programs, and the allowlist root. The executor, the just-in-time compiler, trace generation and the GPU kernels bear only on completeness: an error there yields a proof that fails to verify, not one that verifies falsely. For the chips, the determinism theorems show at most one behaviour per input, where an omission would be silent, and the conformance vectors are the evidence that this behaviour is the ISA's. They cover the core chips; the recursion programs are a hypothesis of Theorem~\ref{thm:composition}.

\subsubsection{ISA conformance}
\label{sec:verification-conformance}
\label{sec:verification-isa}

\paragraph{Vectors from the specification.} The encoding tables of the MIPS32 manual were transcribed into a machine-readable table of 77 instructions, reproduced with the chip proving each in Appendix~\ref{sec:app-isa}. For each, a generator emits ten programs with pseudo-random operands, boundary values and the delay-slot placements the instruction admits, assembled with \texttt{llvm-mc}. The reference is the Unicorn/QEMU MIPS32 emulator~\citep{unicorn}, which shares no code with \Ziren{}'s; each vector compares the general registers, \texttt{HI}/\texttt{LO} and touched memory. All 770 vectors agree, on both the interpreter and the JIT.

\paragraph{Shared instruction suites.} Fifteen hand-written suites, 144 cases, cover instructions with difficult semantics: signed and unsigned division, multiply-accumulate, bit-field and byte-order instructions, rotates, and leading-one and leading-zero counts. One crate supplies case data to both emulator and prover tests, as it does for the generated vectors, so a regression in either appears as a disagreement.

\paragraph{The Cannon suite.} Optimism's Cannon ships hand-written MIPS test programs~\citep{cannon} from the fault-proof project whose semantics a hybrid fault and validity proof must match, exercising the machine against another team's reading of the manual. A harness byte-swaps them to little-endian and adapts the halt convention; all 49 programs applicable to a little-endian guest pass.

\paragraph{From vectors to proofs.} 224 of the generated programs and 48 of the 49 Cannon programs were also proved and verified end to end. The exception halts with a non-zero exit code, which the emulator accepts and the arithmetisation does not admit. Elsewhere the arithmetisation accepts every behaviour the reference exhibits. The next subsection gives part of the converse: the arithmetisation admits at most one behaviour per input; that this behaviour is the ISA's rests on the vectors above.

\paragraph{Executable formal model.} A \Lean{} model of the ISA was written independently of \Ziren{}'s emulator: a decoder transcribed from the encoding tables, an interpreter with the delay-slot program-counter pair and little-endian byte memory, and an ALU. Two checks, both by \texttt{native\_decide}, connect it to the implementation: its decoder agrees with the emulator's on all 2\,275 distinct instruction words in the vectors, and it reproduces the reference results on all 770 vectors. Digests of the emitted event records also agree between repeated runs and between interpreter and JIT, which is evidence of reproducibility rather than a proof of determinism.

\subsubsection{Determinism of the chips}
\label{sec:verification-det}

\paragraph{Statement.} The statement for a chip $U$ is Definition~\ref{def:det} for its constraint system $C_U$, with byte-table lookups lowered to range constraints or to an abstract table relation $T$. A theorem that uses such a relation takes $\forall i, o, o'.\ T(i,o) \wedge T(i,o') \Rightarrow o = o'$ as a hypothesis; 32 of the \ndettotal{} do, for the AND, XOR, OR, NOR and shift-carry operations of the \texttt{Byte} table, which is a function because the verifying key fixes it. Each theorem also takes the equality of a list of values assumed deterministic, which is empty for every chip. The inputs $I_U$ include the declared free values, which Table~\ref{tab:det-findings} records. These theorems are the premises of Theorem~\ref{thm:unique}.

Reading a lookup as a fact about the row that sends it is sound only under a side condition on the whole trace: the lookup argument equates multisets, so a tuple sent with multiplicity $+1$ by one row and $-1$ by another needs no table entry, and every multiplicity must therefore be non-negative on every row, including the padding rows that fill a chip's height. Two further statements per chip discharge it. With $\mathsf{is\_real}$ and every selector at zero, every lookup multiplicity of the row is zero, so a padding row takes part in no bus and no table; and on a real row every multiplicity is a bit. Both are checked by the propagation analysis of the next paragraph and emitted as \Lean{} theorems beside the determinism statements. A chip that leaves a multiplicity unconstrained off its real rows fails the first, which is how the \texttt{SysLinux} gap of Table~\ref{tab:det-findings} was found.

\paragraph{Extraction.} The Rust constraint description that the prover evaluates and the verifier checks is evaluated once with symbolic columns; the resulting constraints, range facts and bus interactions are lowered to a first-order statement, and a generator emits one \Lean{} file per chip with a row record, the constraints, the input and output projections, and the theorem. Selector-specialised chips get one theorem per opcode selector, all restricted to real rows. The extraction covers every core chip of Appendix~\ref{sec:app-chips} except the three preprocessed tables, which carry no statement, and the two SHA-256 control chips, which it currently models as opaque helpers: \ndettotal{} theorems over 57 chips. Nothing in the pipeline is specific to a chip: modelling those two, or the recursion machine's chips, which are written against the same interface, extends it without a new technique.

\paragraph{Proof automation.} As in Picus~\citep{picus}, a propagation analysis over the extracted constraints derives every output column of $w'$ from the inputs as a sequence of steps: linear and inverse solving, bit and one-hot decoding, integer bounds, carry roots, and \emph{gadget summaries} for recurring sub-circuits (a field operation with its carry rows, a byte comparison, a leading-one detector, a canonical word), with case splits on selector bits both witnesses share. The generator replays the derivation in \Lean{}. Each step becomes a lemma stated over only the columns and conjuncts it uses, closed by a small tactic that lifts its statement from $\F$ to the integers with each column's range and decides it with \texttt{omega}, \texttt{linear\_combination} or \texttt{grind}. Not everything is generated: the gadget library is written by hand, the accelerator, division and \texttt{Global} gadget proofs come from per-family generator scripts, and the selector-flag postconditions of ten chips are proved by hand where the automation runs out of budget.

\paragraph{Large chips.} The gadget summaries are library theorems proved once; an accelerator whose outputs come from chained field operations or permutation rounds gets a generated proof that its output columns agree (\texttt{gadget\_det}), split into a chain of segment lemmas of a few hundred steps, each taking the facts live at its start and ending by applying the next. The chip's theorem never splits $w$ and $w'$ into their thousands of columns: it applies the step lemmas to projections of the two constraint sets, each step lemma restating its conjuncts in their original form and substituting what earlier steps proved inside its own small context. A theorem elaborated as one declaration in the full context of a precompile chip took hours; in this form every module of the check takes minutes.

\paragraph{Findings.} A determinism statement fails exactly when a row admits two witnesses with the same inputs and different outputs, so every statement the automation could not close was read for the witness that breaks it. Table~\ref{tab:det-findings} lists what that reading found, together with what the padding statement found. Four are soundness gaps in the constraints, three exploitable and one (\texttt{Global}) latent while every sender sets a flag, each closed by new constraints; the others are prover freedom the protocol allows by design, which the statements now name as inputs rather than hide.

\begin{table}[t]
\centering
\small
\caption{What the determinism statements exposed. Each row is a chip whose statement, as first extracted, did not hold; the last column is how the constraints or the statement were changed. The four new constraint sets change the verifying key.}
\label{tab:det-findings}
\begin{tabular}{@{}>{\raggedright\arraybackslash}p{2.8cm}>{\raggedright\arraybackslash}p{5.4cm}>{\raggedright\arraybackslash}p{3.4cm}>{\raggedright\arraybackslash}p{3.6cm}@{}}
\toprule
Chip & What the statement exposed & Consequence & Resolution \\
\midrule
\texttt{*DoubleAssign} (secp256k1, secp256r1, BN254, BLS12-381) &
  The slope gadget proves $s \cdot 2y = 3x^2 + a$ and nothing about $2y$. For an input with $y \equiv 0$, such as $(0,0)$ on a curve with $a = 0$, every slope satisfies the chip, and the row writes $(s^2, -s^3)$. &
  Soundness gap: the written point is the prover's choice. An AIR-level test forges it for three slopes. &
  An inverse check proves $\mathsf{is\_real}/2y$ exists, and the executor refuses $y \equiv 0$. \\
\addlinespace
\texttt{Global} (flags) &
  \texttt{is\_send} and \texttt{is\_receive} choose the sign of the lifted point's $y$ and were neither boolean nor tied to \texttt{is\_real}: with both zero, $(x, y)$ and $(x, -y)$ both satisfy the row and give different digests. &
  Soundness rests on the senders: every current sender sets exactly one flag, so a $(0,0)$ row cannot balance the bus today. SP1 v6 constrains the flags the same way. &
  Both flags boolean and $\mathsf{is\_send} + \mathsf{is\_receive} = \mathsf{is\_real}$. \\
\addlinespace
\texttt{KeccakSponge\-Control} &
  Of a multi-block sponge, only the first block's input and output addresses were bound to the syscall; later blocks' addresses were free, and nothing tied the number of blocks to the input length. &
  Soundness gap: a prover can absorb later blocks from any address, stop early or late, and write the digest anywhere. &
  Each block carries the input address, output address and length on the block bus; the first block is block 0 and binds the length read from memory, its top byte range-checked; the final block has $(b+1) \cdot 36 = \mathsf{len}$ (36 words, the chip's 1152-bit rate). \\
\addlinespace
\texttt{SysLinux} (padding rows, this work) &
  The gadget that proves $\mathsf{v_0} = \max(a_0, \mathsf{brk})$ sends its complementary byte lookup with a multiplicity constrained only on \texttt{brk} rows, and the decoded call flags that gate the chip's memory accesses were constrained only on real rows; a padding row could carry any of them, including $-1$. &
  Soundness gap: a padding row cancels the lookups of a \texttt{brk} row that lies about the comparison, so the call returns the old break; a padding row can also move the heap register without a call. &
  The multiplicity is the flag sum times the gate, the decode results vanish off real rows, $\mathsf{is\_real}$ is boolean and the recognised calls sum to it. \\
\addlinespace
\texttt{Global} &
  The \texttt{lift\_x} offset is a byte the prover picks among those that lift the message to a curve point. &
  None: the encoding stays injective because the top message field is below $2^{16}$; SP1 is the same. &
  The offset is an input of the statement. \\
\addlinespace
\texttt{SyscallInstrs} &
  The length of a hint read is chosen by the host. &
  None: hints are untrusted input by definition. &
  The hint length is an input of the statement. \\
\addlinespace
\texttt{MemoryGlobalInit} &
  The address and value of an initial-memory row are the prover's choice. &
  None: the global memory argument fixes them against the image and the first access. &
  Address and value are inputs of the statement. \\
\bottomrule
\end{tabular}
\end{table}

{\sloppy\paragraph{Coverage.} The multiplicity statements cover the same 62 chips. In 35 of them every multiplicity is a product with $\mathsf{is\_real}$ or a selector, so the padding statement holds by substitution; 23 carry a padding statement with content, 328 multiplicities in all, of which the analysis proves 286 zero, and the 42 it cannot are the five columns of \texttt{SysLinux} in Table~\ref{tab:det-findings}, all proved once the gates of its resolution are in place; the three preprocessed tables and the shard-boundary memory chip have no row the environment can mark as padding. On real rows 608 of 637 multiplicities are proved bits, 28 more are bits given that an opcode flag the row receives from the program table is one, and the one open, again \texttt{SysLinux}'s comparison gate, closes with the same fix. The determinism statements: all \ndettotal{} theorems are closed: 23 for the integer ALU chips, 8 for the shifts, 9 for control flow, 14 for the memory instructions, 5 for the memory argument, 15 for the miscellaneous and syscall chips, and 30 for the accelerators, one per chip or per opcode-selector specialisation of it. For each, \texttt{\#print axioms} lists only \texttt{propext}, \texttt{Classical.choice} and \texttt{Quot.sound}, so no proof rests on \texttt{sorry} or an added axiom, and every one of the 7\,414 modules of the 62 chip files (57 with statements) builds without a \texttt{sorry}. The check splits each chip file into modules by dependency and builds them in parallel on the evaluation host of \S\ref{sec:perf-setup}; with a 30-minute cap per module, the longest takes 22 minutes. The count is for the analysed constraint system. The checked chip files, gadget snippets and check scripts are archived\footnote{\url{https://zkm-toolchain.s3.amazonaws.com/fv/ziren-fv-lean-20261001b.tar.gz}.} with a per-chip table of theorems, modules and module times; the archived generators regenerate them byte for byte.\par}


\subsubsection{A negative control for the height claims}
\label{sec:verification-forgery}

The verifying key depends on the chip set (\S\ref{sec:recursion}) and the committed geometry is quantised from declared heights (\S\ref{sec:pcs-geometry}), so the argument must reject a shard whose declared shape is false. A harness proves real shard proofs for an arithmetic and a hash-heavy guest and tampers with them in eighteen ways: over- and under-claimed heights, false chip activity, and tampered column and row counts, each injected into the degree mask, the Fiat--Shamir prologue, or both. Every case is rejected: height and activity lies in the degree mask by the \LogUp{} last-layer reconstruction, lies confined to the transcript by the grinding witness (the prologue observes the declared heights), a tampered column count by the jagged verifier, and a tampered row count by the preprocessed round, which reads it from the verifying key.
